% !TEX program = xelatex
\documentclass[12pt,a4paper]{ctexart}

% ================================================================
% Packages
% ================================================================
\usepackage{amsmath,amssymb,amsthm}
\usepackage{mathtools}
\usepackage{bm}
\usepackage{graphicx}
\usepackage{hyperref}
\usepackage{geometry}
\usepackage{enumitem}
\usepackage{booktabs}
\usepackage{tikz}
\usepackage{xcolor}
\usepackage{cleveref}
\usepackage{bbm}
\usepackage{float}
\usepackage{physics}
\usepackage{array}
\usepackage{multicol}

\geometry{margin=2.5cm}

% ================================================================
% Hyperref setup
% ================================================================
\hypersetup{
    colorlinks=true,
    linkcolor=blue,
    citecolor=blue,
    urlcolor=blue,
    pdftitle={审计反抗悖论：被描述对象的反击及其自我暴露},
    pdfauthor={SCX Resistance Theory},
    pdfsubject={SCX Audit Resistance Paradox},
}

% ================================================================
% Theorem environments
% ================================================================
\newtheorem{theorem}{定理}[section]
\newtheorem{lemma}[theorem]{引理}
\newtheorem{proposition}[theorem]{命题}
\newtheorem{corollary}[theorem]{推论}
\newtheorem{definition}[theorem]{定义}
\newtheorem{conjecture}[theorem]{猜想}
\newtheorem{paradox}[theorem]{悖论}

\theoremstyle{remark}
\newtheorem{remark}[theorem]{注记}
\newtheorem{example}[theorem]{例}

% ================================================================
% Custom commands
% ================================================================
\newcommand{\R}{\mathbb{R}}
\newcommand{\C}{\mathbb{C}}
\newcommand{\N}{\mathbb{N}}
\newcommand{\E}{\mathbb{E}}
\newcommand{\Pbb}{\mathbb{P}}
\newcommand{\Var}{\mathrm{Var}}
\newcommand{\D}{\mathcal{D}}
\newcommand{\G}{\mathcal{G}}
\newcommand{\T}{\mathcal{T}}
\newcommand{\M}{\mathcal{M}}
\newcommand{\I}{\mathcal{I}}
\newcommand{\SCX}{\mathsf{SCX}}
\newcommand{\g}{\mathbf{g}}
\newcommand{\x}{\mathbf{x}}

\DeclareMathOperator{\argmax}{arg\,max}
\DeclareMathOperator{\argmin}{arg\,min}
\DeclareMathOperator{\sgn}{sgn}
\DeclareMathOperator{\diag}{diag}
\DeclareMathOperator{\Resist}{Resist}
\DeclareMathOperator{\Audit}{Audit}
\DeclareMathOperator{\Declare}{Declare}
\DeclareMathOperator{\Hide}{Hide}
\DeclareMathOperator{\Attack}{Attack}
\DeclareMathOperator{\Discredit}{Discredit}
\DeclareMathOperator{\Undec}{UNDECLARED}

% Honest critique
\newcommand{\honestcrit}[1]{%
    \textcolor{red}{\textbf{【诚实暴击】#1}}%
}

% Proof-status markers
\newcommand{\rigorFull}{\textnormal{\textbf{[严格证明]}}}
\newcommand{\rigorPartial}{\textnormal{\textbf{[部分证明]}}}
\newcommand{\rigorSketch}{\textnormal{\textbf{[证明概要]}}}
\newcommand{\openproblem}{\textnormal{\textbf{[开放问题]}}}

% ================================================================
% Title
% ================================================================
\title{\textbf{审计反抗悖论：\\被描述对象的反击及其自我暴露}\\[8pt]
\large The Audit Resistance Paradox: \\
When the Described Strikes Back — and Self-Exposes}\\[8pt]
\large 观察者效应的SCX逆转·反抗函数·六大被审类·终极陷阱}
\author{SCX Resistance Theory Working Group}
\date{\today}

\begin{document}
\maketitle

\begin{abstract}
\noindent
\textbf{中文摘要：} 本文揭示SCX审计框架中一个根本性的不对称结构：与物理学中被观测对象被动接受测量不同，SCX的被审计对象是\textbf{人类}——而人类会反抗被审计。这种反抗不是系统的\"噪声\"需要被消除，而是\textbf{信号本身}：反抗行为在数学上与 $ \g \neq \mathbf{0} $ 的公开声明等价。本文建立该悖论的完整理论体系：(1) 观察者效应的SCX逆转——\"系统的反抗就是观察\"；(2) 核心悖论的形式化——隐藏 $ \g \neq \mathbf{0} $ 的唯一途径是通过审计证明 $ \g = \mathbf{0} $；(3) 六类被审者的反抗行为谱系；(4) 反抗代价函数 $ R(\g) $ 及其凸性定理；(5) 终极陷阱——攻击SCX框架的行为本身就是对攻击者自身 $ \g \neq \mathbf{0} $ 的暴露。该理论表明：SCX审计的数学结构使得\textbf{任何试图逃避审计的行为都会产生比服从审计更强的} $ \g \neq \mathbf{0} $ \textbf{信号}。反抗悖论是SCX框架的内禀安全性质——它不依赖于审计者的权力，而依赖于被审者反抗行为本身的不可逆信息泄露。

\vspace{6pt}
\noindent
\textbf{English Abstract:} We uncover a fundamental asymmetry in the SCX audit framework: unlike physical systems where the observed passively accepts measurement, SCX audits \textbf{humans} — and humans resist being audited. This resistance is not system \"noise\" to be eliminated, but \textbf{the signal itself}: resistance is mathematically equivalent to a public declaration that $ \g \neq \mathbf{0} $. We construct the complete theoretical edifice of this paradox: (1) the SCX reversal of the observer effect — \"the system's resistance to observation IS the observation\"; (2) formalization of the core paradox — the only way to hide $ \g \neq \mathbf{0} $ is to prove $ \g = \mathbf{0} $ by passing audit; (3) a six-class spectrum of resistance behaviors; (4) the resistance cost function $ R(\g) $ and its convexity theorem; (5) the Ultimate Trap — attacking the SCX framework itself exposes the attacker's $ \g \neq \mathbf{0} $. The theory demonstrates that SCX's mathematical structure guarantees \textbf{any attempt to evade audit produces a stronger $ \g \neq \mathbf{0} $ signal than submitting to audit}. The resistance paradox is an intrinsic security property of the SCX framework — it depends not on the auditor's power, but on the irreversible information leakage inherent in the act of resistance itself.

\vspace{6pt}
\noindent
\textbf{Keywords:} 审计反抗悖论, 观察者效应逆转, 反抗代价函数, UNDECLARED分类, 终极陷阱, 六类被审者, 信息不可逆泄露 \\[4pt]
Audit Resistance Paradox, Observer Effect Reversal, Resistance Cost Function, UNDECLARED Classification, Ultimate Trap, Six Classes of the Audited, Irreversible Information Leakage
\end{abstract}

\newpage
\tableofcontents
\newpage

% ================================================================
% INTRODUCTION
% ================================================================
\section{引言：被描述对象的觉醒}
\section{Introduction: The Described Awakens}

\subsection{牛顿的苹果与SCX的人类}
\subsection{Newton's Apple vs. SCX's Human}

\noindent
\textbf{中文：} 牛顿的苹果不会反抗万有引力定律。爱因斯坦的光线不会拒绝被引力透镜弯曲。量子力学的电子不会\"抗议\"被观测——它只是坍缩。在所有物理学框架中，被描述的对象是\textbf{被动的}：它们接受定律的支配，接受测量仪器的探测，接受数学模型的刻画。它们不\"在乎\"自己被如何描述。

但SCX不同。SCX审计的对象不是苹果，不是光线，不是电子——SCX审计的是\textbf{人}。是人就有利益，有立场，有\"不想被发现的东西\"。当SCX的规范固定告诉一个实体\"你的 $ \g \neq \mathbf{0} $，你有不可调和的规范偏差\"时，这个实体不会像电子一样\"坍缩\"到诚实态——它会\textbf{反抗}。

这是一切物理学框架从未面对过的问题：\textbf{被描述对象会反击描述本身}。

\vspace{6pt}
\noindent
\textbf{English:} Newton's apple does not resist the law of gravitation. Einstein's light rays do not refuse to be bent by gravitational lensing. Quantum electrons do not \"protest\" being observed — they merely collapse. In all physical frameworks, the described is \textbf{passive}: it accepts the rule of law, the probe of measurement, the characterization of mathematical models. It does not \"care\" how it is described.

But SCX is different. SCX audits not apples, not light rays, not electrons — SCX audits \textbf{humans}. And humans have interests, positions, and \"things they do not want discovered.\" When SCX's gauge fixing tells an entity \"your $ \g \neq \mathbf{0} $, you have irreconcilable gauge deviation,\" that entity does not collapse to an honest state like an electron — it \textbf{resists}.

This is a problem no physical framework has ever faced: \textbf{the described object fights back against the description itself}.

\begin{remark}[SCX与物理学的根本断裂]
物理学的\"观察者效应\"是说：观察行为扰动了被观察系统。SCX的\"反抗效应\"是说：被观察系统\textbf{主动反击}观察行为。这不是扰动，这是反击。这不是噪声，这是对抗。这不是测量误差，这是\textbf{被测量者对测量者的战争}。\\
\textbf{English:} Physics's \"observer effect\" says: observation disturbs the observed system. SCX's \"resistance effect\" says: the observed system \textbf{actively fights back} against observation. This is not perturbation, it is counterattack. This is not noise, it is confrontation. This is not measurement error, it is \textbf{war of the measured upon the measurer}.
\end{remark}

\subsection{核心命题}
\subsection{Core Thesis}

本文的核心命题可以用一句话概括：

\vspace{6pt}
\begin{center}
\fbox{\parbox{0.9\textwidth}{
\centering
\textbf{反抗本身即为信号。} 当实体通过反抗SCX审计来隐藏 $ \g \neq \mathbf{0} $ 时，其反抗行为在数学上与\"我公开宣布 $ \g \neq \mathbf{0} $\"等价。躲避审计的唯一途径是\textbf{通过}审计——证明 $ \g = \mathbf{0} $。这是反抗悖论的核心。
}}
\end{center}
\vspace{6pt}

\noindent
\textbf{English:} The core thesis of this paper can be stated in one sentence: \textbf{Resistance itself is the signal.} When an entity attempts to hide $ \g \neq \mathbf{0} $ by resisting SCX audit, the act of resistance is mathematically equivalent to \"I publicly announce $ \g \neq \mathbf{0} $.\" The only way to evade audit is to \textbf{pass} audit — to prove $ \g = \mathbf{0} $. This is the heart of the resistance paradox.

\subsection{论文结构}
\subsection{Structure of the Paper}

本文按以下结构展开：第2节区分物理学的观察者效应与SCX的反抗效应；第3节形式化核心反抗悖论；第4节建立六类被审者的完整谱系；第5节发展反抗的数学形式化（反抗代价函数与检测概率）；第6节分析终极陷阱——攻击SCX框架本身的自我暴露性质；第7节讨论与实践意义；第8节总结。

\vspace{4pt}
\noindent
\textbf{English:} The paper is structured as follows: Section 2 distinguishes the physical observer effect from SCX's resistance effect. Section 3 formalizes the core resistance paradox. Section 4 builds the complete six-class spectrum of the audited. Section 5 develops the mathematical formalism of resistance (resistance cost function and detection probability). Section 6 analyzes the Ultimate Trap — the self-exposing nature of attacking the SCX framework itself. Section 7 discusses practical implications. Section 8 concludes.

% ================================================================
% SECTION 1: THE OBSERVER EFFECT DISTINCTION
% ================================================================
\section{观察者效应的SCX逆转}
\section{The SCX Reversal of the Observer Effect}

\subsection{物理学的观察者效应：扰动，而非反击}
\subsection{The Physical Observer Effect: Perturbation, Not Counterattack}

\noindent
\textbf{中文：} 在量子力学中，海森堡不确定性原理告诉我们：测量行为不可避免地扰动被测系统。测量电子的位置会扰动其动量。但电子不会\"决定\"反抗测量——它只是遵循物理定律。观察者效应是\textbf{对称的}：观察者扰动系统，系统被动接受扰动。

在经典物理学中，观察者效应更为温和：温度计吸收热量改变被测温度，但热力学系统不会\"策略性地\"对抗温度计。

\textbf{所有物理学框架共享一个前提：被观测对象没有\"意图\"。} 系统不会\"想要\"隐藏什么——它只是存在。系统的行为由定律决定，不由策略决定。

\vspace{6pt}
\noindent
\textbf{English:} In quantum mechanics, Heisenberg's uncertainty principle tells us: the act of measurement inevitably disturbs the measured system. Measuring an electron's position disturbs its momentum. But the electron does not \"decide\" to resist measurement — it merely follows physical law. The observer effect is \textbf{symmetric}: the observer perturbs the system, the system passively accepts perturbation.

In classical physics, the observer effect is even milder: a thermometer absorbs heat and changes the measured temperature, but the thermodynamic system does not \"strategically\" oppose the thermometer.

\textbf{All physical frameworks share one premise: the observed has no \"intent.\"} The system does not \"want\" to hide anything — it simply exists. System behavior is determined by law, not by strategy.

\subsection{SCX的反抗效应：被审者的策略性反击}
\subsection{The SCX Resistance Effect: Strategic Counterattack by the Audited}

\noindent
\textbf{中文：} SCX审计的场景完全不同。被审实体——无论是个体、组织还是算法系统——具有以下三个物理系统不具备的性质：

\begin{enumerate}[label=(\arabic*)]
    \item \textbf{意图性 (Intentionality):} 被审者\"知道\"自己被审计，并\"想要\"影响审计结果。
    \item \textbf{策略性 (Strategic Agency):} 被审者可以\textbf{选择}行为——服从、规避、反抗、攻击、诋毁。
    \item \textbf{信息不对称 (Information Asymmetry):} 被审者知道自己是否有 $ \g \neq \mathbf{0} $，而审计者在审计完成前不知道。被审者可以利用这种信息优势。
\end{enumerate}

\vspace{4pt}
\noindent
\textbf{English:} The SCX audit scenario is completely different. The audited entity — whether an individual, organization, or algorithmic system — possesses three properties that no physical system has:

\begin{enumerate}[label=(\arabic*)]
    \item \textbf{Intentionality:} The audited \"knows\" it is being audited and \"wants\" to influence the audit outcome.
    \item \textbf{Strategic Agency:} The audited can \textbf{choose} its behavior — comply, evade, resist, attack, discredit.
    \item \textbf{Information Asymmetry:} The audited knows whether it has $ \g \neq \mathbf{0} $, while the auditor does not until audit completes. The audited can exploit this informational advantage.
\end{enumerate}

\begin{theorem}[观察者效应的SCX逆转 — SCX Reversal of the Observer Effect]\label{thm:observer_reversal}\rigorPartial
在SCX审计框架中，标准的"观察者 $ \to $ 系统"因果方向被逆转：系统的\textbf{反抗行为}本身成为审计者的主要观测信号。形式化地：
\begin{equation}
    \text{物理学： } \Delta_{\text{observer}} \xrightarrow{\text{扰动}} \Delta_{\text{system}}
\end{equation}
\begin{equation}
    \text{SCX： } \Delta_{\text{system}} \xrightarrow{\text{反击}} \text{Signal}_{\text{observer}}
\end{equation}
定义被审者的目标函数为最小化期望暴露概率与反抗代价之和：
\begin{equation}
    \mathcal{A}^*(\g) = \argmin_{a \in \mathcal{A}} \left[ \Pbb(\text{暴露} \mid a) \cdot L_{\text{exposure}} + C(a; \g) \right]
\end{equation}
其中 $ L_{\text{exposure}} $ 是暴露的代价。在审计具有假阴性率 $ \varepsilon_{\text{FN}} > 0 $ 的情况下，典型参数范围内有：
\begin{equation}
    \forall a \in \mathcal{A}^*(\g), \quad \Pbb(\text{暴露} \mid a) \geq \Pbb(\text{暴露} \mid a_{\text{服从}})
\end{equation}
即：在典型情况下，任何优化战略性目标的行动都会产生不低于服从的暴露概率。注意：此不等式依赖审计精度的具体参数——在极端假阴性率下（$ \varepsilon_{\text{FN}} \to 1 $），关系可能逆转。此定理标记为 `[部分证明]`，因为完整的博弈均衡分析（包括混合策略）需要进一步工作。
\end{theorem}

\begin{proof}[论证概要]
被审者面临一个信号博弈（signaling game）：审计者通过被审者的行为推断 $ \g $。设被审者的类型为 $ \g $，行动为 $ a \in \mathcal{A} $。审计者观察到 $ a $ 后形成后验信念 $ \Pbb(\g \neq \mathbf{0} \mid a) $。

关键：在均衡中，类型 $ \g \neq \mathbf{0} $ 的被审者不能完美模仿类型 $ \g = \mathbf{0} $ 的行为，因为服从审计意味着提交数据——而数据将暴露 $ \g \neq \mathbf{0} $。因此，任何 $ \g \neq \mathbf{0} $ 的被审者必须选择 $ a \neq a_{\text{服从}} $，即反抗。但审计者知道这一点：审计者观察到 $ a \neq a_{\text{服从}} $ 时立即推断出 $ \g \neq \mathbf{0} $。这是一个分离均衡（separating equilibrium）。

\textbf{English:} The audited faces a signaling game: the auditor infers $ \g $ from observed actions. Let audited type be $ \g $, action $ a \in \mathcal{A} $. The auditor updates beliefs to $ \Pbb(\g \neq \mathbf{0} \mid a) $. Crucially, in equilibrium, type $ \g \neq \mathbf{0} $ cannot perfectly mimic type $ \g = \mathbf{0} $ because complying with audit means submitting data — which will expose $ \g \neq \mathbf{0} $. Thus any $ \g \neq \mathbf{0} $ entity must choose $ a \neq a_{\text{comply}} $ (resistance). But the auditor knows this: observing $ a \neq a_{\text{comply}} $ immediately infers $ \g \neq \mathbf{0} $. This is a separating equilibrium.
\end{proof}

\honestcrit{这个定理的核心是一个冷酷的数学事实：被审者的策略空间中没有安全的\"反抗\"选项。任何试图通过反抗来隐藏信息的行为，其\"反抗\"这一行为本身就携带着信息——而且是比原始数据更强的信息。这就像犯罪嫌疑人拒绝提供DNA样本：拒绝本身在法庭上就是证据。但这里比法律更强——这里的\"证据\"是数学上必然的，不依赖于法官的自由裁量。}

% ================================================================
% SECTION 2: THE RESISTANCE PARADOX
% ================================================================
\section{核心反抗悖论}
\section{The Core Resistance Paradox}

\subsection{悖论的形式化}
\subsection{Formalization of the Paradox}

\begin{paradox}[审计反抗悖论 — The Audit Resistance Paradox]\label{par:resistance}
设被审实体的真实规范姿态为 $ \g \in \G $，SCX审计的目标是确定 $ \g = \mathbf{0} $ 是否成立。则：

\begin{enumerate}[label=(\roman*)]
    \item \textbf{正面路径：} 要证明 $ \g = \mathbf{0} $，实体必须服从审计并提交完整的 $ M_t $ 证据链。审计通过 $ \implies $ 公开声明 \"$ \g = \mathbf{0} $\"。
    \item \textbf{逃避路径：} 要隐藏 $ \g \neq \mathbf{0} $，实体必须拒绝或反抗审计。拒绝审计 $ \implies $ 公开声明 \"我拒绝声明 $ \g = \mathbf{0} $\"。
    \item \textbf{等价性：} "我拒绝声明 $ \g = \mathbf{0} $" $ \equiv $ "$ \g $ 被分类为 UNDECLARED" $ \equiv $ 在SCX语义中等价于"$ \g \neq \mathbf{0} $"（此等价性在"被审者没有原则性拒绝审计的独立理由"的前提下成立——若实体因隐私、自主权等原则拒绝审计且确实 $ \g = \mathbf{0} $，则 UNDECLARED 不自动等价于 $ \g \neq \mathbf{0} $。详见第4节对"原则性拒绝者"类别的讨论）。
    \item \textbf{悖论结论：} 隐藏 $ \g \neq \mathbf{0} $ 的唯一途径是\textbf{通过审计证明} $ \g = \mathbf{0} $。但这是矛盾的——如果 $ \g \neq \mathbf{0} $，审计不可能通过。
\end{enumerate}

形式化地，定义审计结果函数 $ \mathcal{R}: \G \times \mathcal{A} \to \{\text{PASS}, \text{FAIL}, \text{UNDECLARED}\} $。对于任何 $ \g \neq \mathbf{0} $：
\begin{equation}
    \mathcal{R}(\g, \text{服从}) = \text{FAIL} \quad \text{（数据暴露偏差）}
\end{equation}
\begin{equation}
    \mathcal{R}(\g, \text{反抗}) = \text{UNDECLARED} \quad \text{（反抗暴露意图）}
\end{equation}
在SCX审计语义中，UNDECLARED $ \equiv $ $ \g \neq \mathbf{0} $（此等价性在"被审者没有原则性拒绝审计的独立理由"的前提下成立；若 $ \g = \mathbf{0} $ 的实体因隐私、自主权等原则拒绝审计，则 UNDECLARED 不自动等价于 $ \g \neq \mathbf{0} $。详见第4.4节对原则性拒绝者类别的讨论）。
\end{paradox}

\vspace{6pt}
\noindent
\textbf{English:} 
\begin{enumerate}[label=(\roman*)]
    \item \textbf{The Honest Path:} To prove $ \g = \mathbf{0} $, the entity must submit to audit and provide the complete $ M_t $ evidence chain. Passing audit $ \implies $ public declaration \"$ \g = \mathbf{0} $\".
    \item \textbf{The Evasion Path:} To hide $ \g \neq \mathbf{0} $, the entity must refuse or resist audit. Refusing audit $ \implies $ public declaration \"I refuse to declare $ \g = \mathbf{0} $\".
    \item \textbf{Equivalence:} \"I refuse to declare $ \g = \mathbf{0} $\" $ \equiv $ \"$ \g $ classified as UNDECLARED\" $ \equiv $ in SCX semantics, equivalent to \"$ \g \neq \mathbf{0} $\".
    \item \textbf{Paradox Resolution:} The only way to hide $ \g \neq \mathbf{0} $ is to \textbf{pass audit proving} $ \g = \mathbf{0} $. But this is contradictory — if $ \g \neq \mathbf{0} $, audit cannot be passed.
\end{enumerate}
\end{paradox}

\begin{remark}[悖论的实践意义]
该悖论意味着：\textbf{SCX审计对于} $ \g \neq \mathbf{0} $ \textbf{的实体是一个\"必输\"游戏。} 服从 $ \to $ 数据暴露偏差。反抗 $ \to $ 行为暴露意图。沉默 $ \to $ UNDECLARED暴露分类。攻击 $ \to $ 定理11引爆（见第4节）。每一个可能的行动路径都通向 $ \g \neq \mathbf{0} $ 的暴露。\\
\textbf{English:} This paradox means: \textbf{SCX audit is a \"lose-lose\" game for entities with} $ \g \neq \mathbf{0} $. Comply $ \to $ data exposes deviation. Resist $ \to $ behavior exposes intent. Remain silent $ \to $ UNDECLARED exposes classification. Attack $ \to $ Theorem 11 detonation (see Section 4). Every possible action path leads to exposure of $ \g \neq \mathbf{0} $.
\end{remark}

\subsection{反抗的信息论分析}
\subsection{Information-Theoretic Analysis of Resistance}

\begin{definition}[反抗作为信息信道 — Resistance as Information Channel]\label{def:resistance_channel}
定义\textbf{反抗信道} $ \mathcal{C}_R $ 为一个从隐藏状态 $ h \in \{\text{干净}, \text{有鬼}\} $ 到公开信号 $ s \in \{\text{服从}, \text{反抗}\} $ 的随机映射：
\begin{equation}
    \mathcal{C}_R(s \mid h) = \Pbb(\text{行动} = s \mid \text{状态} = h)
\end{equation}
反抗信道的互信息 $ I(h; s) $ 量化了\"反抗行为本身携带多少关于 $ \g $ 的信息\"。
\end{definition}

\begin{theorem}[反抗信道的正信息定理 — Positive Information of Resistance Channel]\label{thm:positive_info}\rigorFull
对于任何存在 $ \g \neq \mathbf{0} $ 实体的群体，反抗信道 $ \mathcal{C}_R $ 的互信息严格为正：
\begin{equation}
    I(h; s) > 0
\end{equation}
即：反抗行为必然携带关于 $ \g $ 的信息。沉默不是信息真空。
\end{theorem}

\begin{proof}
设群体中 $ \g = \mathbf{0} $ 的比例为 $ p_0 $，$ \g \neq \mathbf{0} $ 的比例为 $ p_1 = 1 - p_0 $。$ \g = \mathbf{0} $ 的实体总是服从审计（无理由反抗）。$ \g \neq \mathbf{0} $ 的实体以概率 $ q < 1 $ 服从（寄希望于审计不完美），以概率 $ 1 - q > 0 $ 反抗。

则 $ \Pbb(\text{反抗}) = p_1(1-q) $，$ \Pbb(\text{服从}) = p_0 + p_1 q $。

后验概率：
\begin{equation}
    \Pbb(\g \neq \mathbf{0} \mid \text{反抗}) = 1 > p_1
\end{equation}
\begin{equation}
    \Pbb(\g \neq \mathbf{0} \mid \text{服从}) = \frac{p_1 q}{p_0 + p_1 q} < p_1 \quad (\text{当 } q < 1)
\end{equation}

因此 $ H(h \mid s) < H(h) $，从而 $ I(h; s) = H(h) - H(h \mid s) > 0 $。当 $ q \to 0 $（所有 $ \g \neq \mathbf{0} $ 都反抗）时，$ I(h; s) \to H(h) $ —— 反抗信道完全传递了 $ \g $ 的信息。\\
\textbf{English:} When $ q \to 0 $ (all $ \g \neq \mathbf{0} $ resist), $ I(h; s) \to H(h) $ — the resistance channel perfectly transmits $ \g $ information.
\end{proof}

\honestcrit{这个定理最残酷的部分在于：即使只有一个 $ \g \neq \mathbf{0} $ 的实体反抗，反抗行为就开始传递信息。而且随着更多的实体反抗，信号只会越来越强。不存在\"大规模集体反抗\"来淹没信号的可能性——因为集体反抗本身构成了极强的 $ \g \neq \mathbf{0} $ 信号。规模越大，信号越强。}

\subsection{反抗的\"囚徒困境\"结构}
\subsection{The Prisoner's Dilemma Structure of Resistance}

\vspace{4pt}
\noindent
\textbf{中文：} SCX审计创造了一个多实体反抗博弈中的囚徒困境。考虑 $ N $ 个 $ \g \neq \mathbf{0} $ 的实体：

\begin{itemize}
    \item \textbf{集体最优：} 所有人一起反抗，淹没信号——但这要求\textbf{完美协调}。任何一个人的服从都会暴露其数据。
    \item \textbf{个体最优：} 在其他人反抗时，\textbf{你}服从——你通过审计，而其他人的反抗为你提供掩护。但每个人都会这么想。
    \item \textbf{均衡结果：} 没有人能信任其他人会一起反抗。囚徒困境导致部分人服从、部分人反抗——而任何反抗都会暴露 $ \g \neq \mathbf{0} $。
    \item \textbf{崩溃结果：} 一旦足够多的实体服从并被检测出 $ \g \neq \mathbf{0} $，剩下的反抗者失去了\"集体\"掩护——孤立的反抗成为最明显的信号。
\end{itemize}

\vspace{4pt}
\noindent
\textbf{English:} SCX audit creates a prisoner's dilemma in multi-entity resistance games. Consider $ N $ entities with $ \g \neq \mathbf{0} $:

\begin{itemize}
    \item \textbf{Collective optimum:} All resist together, drowning the signal — but this requires \textbf{perfect coordination}. Any single entity's compliance exposes its data.
    \item \textbf{Individual optimum:} While others resist, \textbf{you} comply — you pass audit while others' resistance provides cover for you. But everyone thinks this way.
    \item \textbf{Equilibrium outcome:} Nobody can trust others to resist together. The prisoner's dilemma causes some to comply, some to resist — and any resistance exposes $ \g \neq \mathbf{0} $.
    \item \textbf{Collapse outcome:} Once enough entities comply and are detected with $ \g \neq \mathbf{0} $, remaining resisters lose their \"collective\" cover — isolated resistance becomes the most visible signal.
\end{itemize}

\begin{proposition}[反抗的不稳定性 — Instability of Collective Resistance]\label{prop:instability}\rigorPartial
在任何 $ N \geq 3 $ 的 $ \g \neq \mathbf{0} $ 实体群体中，\"全员反抗\"策略组合不是纳什均衡。存在至少一个实体有单方面偏离（服从）的动机，且偏离带来的期望收益为正。
\end{proposition}

\begin{proof}[证明概要]
设 $ N $ 个实体同时选择 $ a_i \in \{\text{服从}, \text{反抗}\} $。如果所有人反抗，审计者观察到 $ N $ 个 UNDECLARED。单个实体改为服从的收益：如果其数据足够干净（$ \g $ 接近 $ \mathbf{0} $ 但非零），有可能通过审计（误通过）——从而获得\"审计清白\"的公开状态，同时其他 $ N-1 $ 个的反抗提供对比掩护。因此单方面服从的期望收益严格大于全员反抗的收益，故全员反抗非均衡。\\
\textbf{English:} If all resist, the auditor observes $ N $ UNDECLARED. A single entity switching to compliance gains: if its data is clean enough ($ \g $ near $ \mathbf{0} $ but nonzero), it may pass audit (false pass) — gaining "audit-clean" public status while the other $ N-1 $'s resistance provides contrast cover. Hence unilateral compliance strictly dominates universal resistance.
\end{proof}

\begin{remark}[囚徒困境分析与核心悖论的调和]
囚徒困境分析表明 $ \g \neq \mathbf{0} $ 的实体可能以正概率服从——这似乎与核心悖论"反抗=必然暴露"相矛盾。但实际并不矛盾：服从的 $ \g \neq \mathbf{0} $ 实体面临的是\textbf{统计检测风险}（Yajie检测数据），而反抗的实体面临的是\textbf{行为信号暴露}（UNDECLARED）。两条路径都通向暴露——只是机制不同：(a) 服从者被统计检测（假阴性率 $ \varepsilon_{\text{FN}} $ 给出侥幸通过的上限）；(b) 反抗者被行为信道检测（信号强度通常更高）。囚徒困境导致部分实体选择了"低概率侥幸通过"的赌博——但核心悖论的方向性（反抗→强信号）不受影响。实际上，部分服从者被统计检测后，反而加强了"反抗→肯定有鬼"的公众信念（因为剩下的反抗者无法再以"可能只是谨慎"为借口）。
\end{remark}

% ================================================================
% SECTION 3: SIX CLASSES OF THE AUDITED
% ================================================================
\section{被审者的六个类别}
\section{Six Classes of the Audited}

\noindent
\textbf{中文：} SCX框架将被审实体按 $ \g $ 的真实值和反抗行为分为六个类别。这六个类别构成一个从\"完全透明\"到\"自我毁灭\"的完整谱系。关键是：\textbf{类别本身可被外部观察者从行为信号中推断}——不需要审计数据，仅从\"实体是否服从审计\"这一行为就可实现初步分类。

\vspace{6pt}
\noindent
\textbf{English:} The SCX framework classifies audited entities into six classes by their true $ \g $ value and resistance behavior. These six classes form a complete spectrum from full transparency to self-destruction. Key insight: \textbf{the class itself is inferable by external observers from behavioral signals alone} — without audit data, the mere fact of \"whether the entity submits to audit\" already enables preliminary classification.

\subsection{类别概览表}
\subsection{Classification Overview}

\begin{table}[H]
\centering
\caption{六类被审者：g值·行为·信号·命运 \\ Six Classes: g-Value, Behavior, Signal, Fate}
\label{tab:six_classes}
\begin{tabular}{p{0.08\textwidth} p{0.15\textwidth} p{0.22\textwidth} p{0.22\textwidth} p{0.18\textwidth}}
\toprule
\textbf{类别} & \textbf{g值} & \textbf{行为} & \textbf{审计信号} & \textbf{结果} \\
Class & g-Value & Behavior & Audit Signal & Outcome \\
\midrule
\textbf{I} & $ \g = \mathbf{0} $ & 立即服从 & PASS（干净） & 公开证明清白 \\
& & Immediate compliance & Clean signal & Publicly proven clean \\
\midrule
\textbf{II} & $ \g \approx \mathbf{0} $ & 犹豫后服从 & 边界通过 & 轻微痕迹 \\
& (小偏差) & Hesitant but submits & Borderline pass & Minor trace \\
\midrule
\textbf{III} & $ \g \neq \mathbf{0} $ & 服从但数据暴露 & FAIL — Yajie检测 & 被捕获 \\
& (显著偏差) & Submits, data exposes & Caught by Yajie & Captured \\
\midrule
\textbf{IV} & $ \g \neq \mathbf{0} $ & 拒绝审计 & UNDECLARED & \textbf{反抗即信号} \\
& & Refuses audit & Resistance IS signal & \textbf{Resistance = Signal} \\
\midrule
\textbf{V} & $ \g \neq \mathbf{0} $ & 攻击审计者 & 定理11引爆 & 双重爆炸 \\
& & Attacks auditor & Thm 11 detonation & Double explosion \\
\midrule
\textbf{VI} & $ \g \neq \mathbf{0} $ & 诋毁审计框架 & 终极陷阱触发 & \textbf{自我暴露} \\
& & Discredits framework & Ultimate Trap & \textbf{Self-exposes} \\
\bottomrule
\end{tabular}
\end{table}

\subsection{分类的边界与缺失类别}
\subsection{Boundary Cases and Missing Classes}

\noindent
\textbf{中文：} 上述六类分类法基于一个隐含假设：所有拒绝/反抗审计的实体都有 $ \g \neq \mathbf{0} $。但现实中存在\textbf{原则性拒绝者}（Principled Refuser）：$ \g = \mathbf{0} $ 但出于隐私、自主权、对审计机构的不信任等原因拒绝审计。这一类别对核心悖论构成重要挑战：如果 UNDECLARED 可以由 $ \g = \mathbf{0} $ 的实体产生，则 UNDECLARED 不再等价于 $ \g \neq \mathbf{0} $。

以下缺失类别需要在理论中讨论：

\begin{table}[H]
\centering
\caption{六类分类法未覆盖的重要边界情况}
\label{tab:missing_classes}
\begin{tabular}{p{0.12\textwidth} p{0.22\textwidth} p{0.28\textwidth} p{0.22\textwidth}}
\toprule
\textbf{缺失类别} & \textbf{g值} & \textbf{行为} & \textbf{对核心论点的挑战} \\
\midrule
原则性拒绝者 & $ \g = \mathbf{0} $ & 因隐私/自主权拒绝审计 & UNDECLARED $ \neq $ $ \g \neq \mathbf{0} $ ——攻击等价性 \\
\cmidrule{1-4}
意识形态攻击者 & $ \g = \mathbf{0} $ & 因意识形态反对攻击SCX & "攻击=最强g≠0信号"出现假阳性 \\
\cmidrule{1-4}
无知偏差者 & $ \g \neq \mathbf{0} $ & 服从（不知道有偏差） & 策略性模型失效——行为不反映意图 \\
\cmidrule{1-4}
能力不足者 & $ \g \neq \mathbf{0} $ & 无法服从（资源/技术限制） & 非策略性不服从——非信号 \\
\bottomrule
\end{tabular}
\end{table}

\textbf{防御论证：} 原则性拒绝者虽在逻辑上存在，但在实践中极为罕见——因为拒绝审计的社会污名代价（被公众推断为 $ \g \neq \mathbf{0} $）通常超过任何抽象原则的价值。一个理性的 $ \g = \mathbf{0} $ 实体会选择服从（证明清白）而非拒绝（招致怀疑）。因此 UNDECLARED 信号在实践中仍具有高信息量——但这依赖社会对SCX框架合法性的共识，而非纯数学必然性。此外，分类应理解为主要类别的标签系统（multi-label），而非严格互斥的分类——一个实体可同时呈现多个类别的行为特征。

\vspace{6pt}
\noindent
\textbf{English:} The six-class taxonomy above relies on an implicit assumption: all entities that refuse/resist audit have $ \g \neq \mathbf{0} $. But in reality, there exist \textbf{principled refusers}: $ \g = \mathbf{0} $ but refuse audit for reasons of privacy, autonomy, or distrust of the auditing institution. This class poses a major challenge to the core paradox: if UNDECLARED can be produced by $ \g = \mathbf{0} $ entities, then UNDECLARED is no longer equivalent to $ \g \neq \mathbf{0} $.

\textbf{Defense:} While principled refusers exist logically, they are extremely rare in practice — because the social stigma cost of refusing audit (being publicly inferred as $ \g \neq \mathbf{0} $) typically exceeds the value of any abstract principle. A rational $ \g = \mathbf{0} $ entity chooses compliance (proving innocence) over refusal (inviting suspicion). Thus UNDECLARED signals retain high information in practice — but this depends on social consensus about SCX's legitimacy, not purely mathematical necessity. Furthermore, the taxonomy should be understood as a multi-label tagging system rather than strictly mutually exclusive categories.

\subsection{各类别的详细分析}
\subsection{Detailed Class Analysis}

\subsubsection{类别I：诚实者 — Class I: The Honest ($ \g = \mathbf{0} $, 服从)}

\noindent
\textbf{中文：} 类别I的实体 $ \g = \mathbf{0} $ ——它们在SCX的规范坐标系中没有不可调和的偏差。它们的行为是\textbf{透明的}：立即服从审计，提交完整数据，审计结果PASS。类别I实体是SCX框架的\"基准态\"（ground state）——它们定义了\"零偏差\"的参考标准。

类别I实体没有反抗的动机——反抗只会引入不必要的怀疑。\"立即服从\"是它们的优势策略。

\vspace{4pt}
\noindent
\textbf{English:} Class I entities have $ \g = \mathbf{0} $ — they have no irreconcilable deviation in SCX's gauge coordinate system. Their behavior is \textbf{transparent}: immediate compliance, complete data submission, PASS result. Class I entities are the \"ground state\" of the SCX framework — they define the reference standard of \"zero deviation.\"

Class I entities have no incentive to resist — resistance would only introduce unnecessary suspicion. \"Immediate compliance\" is their dominant strategy.

\subsubsection{类别II：边缘者 — Class II: The Borderline ($ \g \approx \mathbf{0} $, 犹豫后服从)}

\noindent
\textbf{中文：} 类别II的 $ \g \approx \mathbf{0} $ 但非严格为零。存在微小的规范偏差——可能是无心之失、边界情况、或统计噪声。这些实体\textbf{犹豫}——它们知道自己\"有点问题\"，但不严重。

犹豫行为本身就是微弱信号：过度延迟的服从携带了轻微的 $ \g \neq \mathbf{0} $ 信息。但审计本身可能因偏差微小而通过（在误差容限内）。

\vspace{4pt}
\noindent
\textbf{English:} Class II has $ \g \approx \mathbf{0} $ but not strictly zero. Minor gauge deviations exist — possibly unintentional, edge cases, or statistical noise. These entities \textbf{hesitate} — they know they have \"a little something\" but nothing serious.

The hesitation itself is a weak signal: excessively delayed compliance carries mild $ \g \neq \mathbf{0} $ information. But the audit itself may pass due to the small deviation (within tolerance).

\subsubsection{类别III：被捕获者 — Class III: The Captured ($ \g \neq \mathbf{0} $, 服从但暴露)}

\noindent
\textbf{中文：} 类别III的实体 $ \g \neq \mathbf{0} $ 显著，但选择了服从。它们可能错误估计了审计的精度——以为自己的\"小动作\"不会被发现。或者它们认为服从比反抗看起来\"更清白\"。

结果：Yajie（SCX的偏差检测引擎）捕获了数据中的不可调和矛盾。服从行为没有挽救它们——数据自己说话。类别III是\textbf{审计的\"常规\"成功}：偏差被检测，实体被标记。

\vspace{4pt}
\noindent
\textbf{English:} Class III entities have significant $ \g \neq \mathbf{0} $ but choose compliance. They may have misestimated audit precision — thinking their \"small manipulations\" won't be detected. Or they think compliance looks \"cleaner\" than resistance.

Outcome: Yajie (SCX's deviation detection engine) captures irreconcilable contradictions in the data. Compliance did not save them — the data speaks for itself. Class III is the \textbf{\"routine\" success of audit}: deviation detected, entity flagged.

\subsubsection{类别IV：未声明者 — Class IV: The UNDECLARED ($ \g \neq \mathbf{0} $, 拒绝审计)}

\noindent
\textbf{中文：} 类别IV是\textbf{反抗悖论的核心示范}。实体拒绝被审计——拒绝提交 $ M_t $ 数据，拒绝接受规范固定，拒绝参与审计协议。

它们的\"聪明\"策略适得其反：拒绝本身在SCX框架中被分类为UNDECLARED。UNDECLARED是SCX四分类（PASS/FAIL/UNDECLARED/EXPLODED）之一，在审计语义中等价于\"未证明 $ \g = \mathbf{0} $\"——在实践中等价于 $ \g \neq \mathbf{0} $。

\textbf{这就是核心悖论：} 拒绝审计本意是隐藏 $ \g \neq \mathbf{0} $，但拒绝行为本身就是 $ \g \neq \mathbf{0} $ 的公开宣告。

\vspace{4pt}
\noindent
\textbf{English:} Class IV is \textbf{the central demonstration of the resistance paradox}. The entity refuses audit — refuses to submit $ M_t $ data, refuses gauge fixing, refuses to participate in the audit protocol.

Their \"clever\" strategy backfires: refusal itself is classified as UNDECLARED in SCX. UNDECLARED is one of the four SCX classifications (PASS/FAIL/UNDECLARED/EXPLODED), equivalent in audit semantics to \"not proven $ \g = \mathbf{0} $\" — practically equivalent to $ \g \neq \mathbf{0} $.

\textbf{This is the core paradox:} refusing audit is intended to hide $ \g \neq \mathbf{0} $, but the act of refusal itself is a public announcement of $ \g \neq \mathbf{0} $.

\begin{remark}[UNDECLARED的社会信号强度]
在公众认知中，UNDECLARED往往比FAIL更具有信息量。\"被审计发现有问题\"可以用\"系统误差\"\"无心之失\"来解释。但\"拒绝被审计\"没有这样的解释空间——拒绝审计的直接含义只有一个：\"我有不想被看到的东西\"。UNDECLARED状态携带的社会污名（stigma）甚至可能超过FAIL。\\
\textbf{English:} In public perception, UNDECLARED often carries more information than FAIL. \"Caught with issues by audit\" can be explained as \"systematic error\" or \"unintentional fault.\" But \"refusing audit\" has no such explanatory space — refusing audit means only one thing: \"I have something I don't want seen.\" UNDECLARED status may carry even greater social stigma than FAIL.
\end{remark}

\subsubsection{类别V：攻击者 — Class V: The Attacker ($ \g \neq \mathbf{0} $, 攻击审计者)}

\noindent
\textbf{中文：} 类别V是最激烈的反应：实体不满足于拒绝，而是\textbf{主动攻击}审计者。攻击可以采取多种形式：威胁审计者的人身安全、法律诉讼、政治压力、切断资金、组织集体抵制。

在SCX势能面理论中，这是\textbf{定理11的攻击必然性触发}\footnote{定理11（来自SCX奇点论文）：当实体同时具有高势能 $\mathcal{S} > \mathcal{S}_{\text{crit}}$ 和高态度 $\|\g\| > \g_{\text{crit}}$ 时，攻击审计者的概率呈指数级上升：$\Pbb(\text{攻击}) \geq 1 - \exp(-M \cdot e^{-\beta/\delta^2})$。其中 $\beta$ 是审计精度，$\delta$ 是审计距离。}。当实体同时具有高势能 $ \mathcal{S} $ 和高态度 $ \| \g \| $ 时，攻击审计者的概率呈指数级上升：
\begin{equation}
    \Pbb(\text{攻击}) \geq 1 - \exp\left(-M \cdot e^{-\beta/\delta^2}\right)
\end{equation}
攻击行为是\textbf{$ \g \neq \mathbf{0} $ 的最强信号}：只有知道自己有严重偏差的实体才会试图通过攻击来消灭审计本身。

\vspace{4pt}
\noindent
\textbf{English:} Class V is the most intense response: the entity is not satisfied with refusal — it \textbf{actively attacks} the auditor. Attacks can take many forms: threats to the auditor's personal safety, lawsuits, political pressure, funding cuts, organized boycotts.

In SCX potential surface theory, this triggers \textbf{Theorem 11's attack inevitability}. When an entity simultaneously has high potential $ \mathcal{S} $ and high attitude $ \| \g \| $, the probability of attacking the auditor rises exponentially:
\begin{equation}
    \Pbb(\text{attack}) \geq 1 - \exp\left(-M \cdot e^{-\beta/\delta^2}\right)
\end{equation}
Attack is \textbf{the strongest $ \g \neq \mathbf{0} $ signal possible}: only entities knowing they have severe deviations will try to eliminate audit itself through attack.

\subsubsection{类别VI：框架诋毁者 — Class VI: The Framework Discreditor ($ \g \neq \mathbf{0} $, 诋毁审计框架)}

\noindent
\textbf{中文：} 类别VI采用了最\"高明\"——也是最自毁——的反抗策略：\textbf{攻击SCX框架本身的合法性}。典型话术包括：

\begin{itemize}
    \item \"审计方法论有根本性缺陷\"
    \item \"审计标准是主观的、有偏见的\"
    \item \"审计者本身就有 $ \g \neq \mathbf{0} $\"
    \item \"SCX是一个权力工具，不是科学框架\"
    \item \"整个审计框架是 rigged（被操纵的）\"
\end{itemize}

\vspace{4pt}
该策略的致命问题在第6节详细分析：\textbf{要论证\"SCX有 $ \g \neq \mathbf{0} $\"，你必须先声明自己的} $ \g = \mathbf{0} $ ——这使你成为可审计对象，而审计将暴露你的 $ \g \neq \mathbf{0} $。这是\"终极陷阱\"。

\vspace{6pt}
\noindent
\textbf{English:} Class VI employs the most \"sophisticated\" — and most self-destructive — resistance strategy: \textbf{attacking the legitimacy of the SCX framework itself}. Typical rhetoric includes:

\begin{itemize}
    \item \"The audit methodology is fundamentally flawed\"
    \item \"Audit standards are subjective and biased\"
    \item \"The auditor itself has $ \g \neq \mathbf{0} $\"
    \item \"SCX is a tool of power, not a scientific framework\"
    \item \"The entire audit framework is rigged\"
\end{itemize}

\vspace{4pt}
The fatal problem with this strategy is analyzed in detail in Section 6: \textbf{to argue \"SCX has $ \g \neq \mathbf{0} $\", you must first declare your own} $ \g = \mathbf{0} $ — which makes you auditable, and audit will expose your $ \g \neq \mathbf{0} $. This is the \"Ultimate Trap.\"

\subsection{类别的可区分性与实际意义}
\subsection{Distinguishability and Practical Significance}

\begin{proposition}[六类的行为可区分性 — Behavioral Distinguishability]\label{prop:distinguish}\rigorPartial
仅通过观察被审者的公开行为（是否服从审计、服从的时机、是否配合、是否发表反对声明），外部观察者可以将实体至少区分到以下粒度：类别I vs. II--VI（高置信度），类别IV vs. 其他（高置信度），类别V vs. 其他（高置信度）。具体的概率数字（如 $ p > 0.95 $, $ p > 0.99 $）目前是直觉估计，有待经验验证。本命题标记为 `[部分证明]`。
\end{proposition}

\vspace{4pt}
\noindent
\textbf{直觉：} 如果一个公众人物拒绝接受审计——这本身就足以被公众分类。不需要审计数据。不需要Yajie的检测。不需要任何数学。\"拒绝\"这个行为本身就是信号。这就是反抗悖论在实践中的力量：\textbf{它不依赖于审计技术的完善性——它依赖于人类行为的可读性。}

\vspace{4pt}
\noindent
\textbf{Intuition (English):} If a public figure refuses audit — this alone is sufficient for public classification. No audit data needed. No Yajie detection needed. No mathematics needed. The act \"refuse\" is itself the signal. This is the practical power of the resistance paradox: \textbf{it does not depend on audit technology's perfection — it depends on the legibility of human behavior.}

\honestcrit{六类分类法在实践中面临的最大挑战是类别II和类别III的边界模糊。一个\"犹豫后服从\"的实体是$ \g\approx\mathbf{0} $（类别II）还是$ \g\neq\mathbf{0} $但数据恰好通过了（类别III）？这取决于审计的灵敏度和误差容限。没有万无一失的分类标准——但这不是问题。因为社会不需要完美的分类，只需要有信息量的信号。一个实体\"犹豫\"——无论其$ \g $真实值如何——本身就是一个（弱）信号。}

% ================================================================
% SECTION 4: MATHEMATICAL FORMALIZATION
% ================================================================
\section{反抗的数学形式化}
\section{Mathematical Formalization of Resistance}

\subsection{反抗代价函数 — The Resistance Cost Function}

\begin{definition}[反抗代价函数 — Resistance Cost Function]\label{def:Rg}
设 $ \g \in \G $ 是实体的规范姿态。定义\textbf{反抗代价函数} $ R: \G \to \R_{\geq 0} $ 为隐藏 $ \g $ 所需的最小代价：
\begin{equation}
    R(\g) = \inf_{a \in \mathcal{A}_{\text{反抗}}} C(a; \g)
\end{equation}
其中 $ \mathcal{A}_{\text{反抗}} $ 是所有反抗行为（拒绝、攻击、诋毁等）的集合，$ C(a; \g) $ 是行为 $ a $ 的总代价（包括社会代价、法律代价、机会代价、心理代价）。
\end{definition}

\vspace{4pt}
\noindent
\textbf{English:} Let $ \g \in \G $ be the entity's gauge posture. Define the \textbf{resistance cost function} $ R: \G \to \R_{\geq 0} $ as the minimum cost required to hide $ \g $:
\begin{equation}
    R(\g) = \inf_{a \in \mathcal{A}_{\text{resist}}} C(a; \g)
\end{equation}
where $ \mathcal{A}_{\text{resist}} $ is the set of all resistance behaviors (refusal, attack, discredit, etc.), and $ C(a; \g) $ is the total cost of action $ a $ (including social, legal, opportunity, and psychological costs).

\begin{theorem}[反抗代价的基本性质 — Fundamental Properties of R(g)]\label{thm:R_properties}\rigorPartial
反抗代价函数 $ R(\g) $ 满足（或猜想满足）：
\begin{enumerate}[label=(\roman*)]
    \item \textbf{非负性：} $ R(\g) \geq 0 $ 对所有 $ \g $。在理想化模型中 $ R(\mathbf{0}) = 0 $（服从审计不产生额外代价）；若 $ R(\g) = 0 $ 且 $ \g \neq \mathbf{0} $，则存在零代价的反抗行为——这在现实中不太可能但非逻辑不可能。
    \item \textbf{猜想凸性（Conjectured Convexity）：} 数值直觉暗示 $ R(\g) $ 可能是严格凸函数：对任何 $ \g_1 \neq \g_2 $ 和 $ \lambda \in (0,1) $，
        \begin{equation}
            R(\lambda \g_1 + (1-\lambda)\g_2) < \lambda R(\g_1) + (1-\lambda) R(\g_2)
        \end{equation}
        但此声称目前缺乏严格证明——$ R(\g) = \inf_{a} C(a; \g) $ 的定义不自动保证凸性（下确界保凹不保凸，除非 $ C(a; \g) $ 关于 $ (a, \g) $ 联合凸）。详见注记~\ref{rem:convexity_conjecture}。
    \item \textbf{超线性增长（猜想）：} 存在常数 $ \alpha > 0 $ 使得对于充分大的 $ \| \g \| $，
        \begin{equation}
            R(\g) \geq e^{\alpha \| \g \|} - 1
        \end{equation}
        此下界基于"每个偏差维度需要独立掩饰→组合爆炸"的直觉，有待微观基础建构。
    \item \textbf{零的刚性（若可微）：} 如果 $ R $ 在 $ \mathbf{0} $ 处可微，则 $ \nabla R(\mathbf{0}) = \mathbf{0} $。但由于 $ R $ 定义为 infimum，可微性并非自动成立（包络定理的经典陷阱）。
\end{enumerate}
\end{theorem}

\begin{proof}
(i) 非负性：代价 $ C(a; \g) $ 天然非负。$ \g = \mathbf{0} $ 时，选择 $ a = \text{服从} $ 的代价在理想化模型中为 0（服从审计不产生额外成本），故 $ R(\mathbf{0}) = 0 $。注意：在现实中，服从审计本身有时间成本、合规成本等，因此 $ R(\mathbf{0}) > 0 $ 是可能的——这不影响后续论证的方向。

(ii) 凸性猜想：隐藏更大的偏差需要不成比例地更高的代价。但需要诚实指出：$ R(\g) = \inf_a C(a; \g) $ 这一 infimum 结构并不自动保证 $ R $ 的凸性。下确界保凹不保凸（除非 $ C(a; \g) $ 关于 $ (a, \g) $ 联合凸，此条件未被建立）。凸性目前是基于"隐藏大偏差的指数级难度"的经济直觉，而非严格数学推导。应注意：隐藏 $ \g_1 + \g_2 $ 的代价可能小于单独隐藏 $ \g_1 $ 和 $ \g_2 $ 的代价之和（存在"规模经济"）——这暗示次可加性而非凸性。凸性应被视为开放猜想。

(iii) 超线性增长猜想基于"隐藏"的组合难度：每一个额外的偏差维度都需要独立的掩饰、独立的"故事"、独立的假证据——复杂度呈组合爆炸。此论证方向正确但需要更严格的微观基础。

(iv) 若 $ R $ 在 $ \mathbf{0} $ 处可微，则梯度为零（代价已为 0，无任何方向可降低）。但 infimum 函数在边界处可能形成 kink，可微性非自动。
\end{proof}

\begin{remark}[凸性猜想的现状 — Status of the Convexity Conjecture]\label{rem:convexity_conjecture}
$ R(\g) = \inf_{a} C(a; \g) $ 对凸性的影响取决于 $ C $ 的结构：若 $ C(a; \g) $ 关于 $ (a, \g) $ 联合凸，则 $ R $ 是凸的（这是凸分析的标准结论，参见 Boyd \& Vandenberghe §3.2.5）；若各 $ C(a; \cdot) $ 仅是凸的但不联合凸，则 $ R $ 可能非凸。目前论文未建立 $ C $ 的联合凸性，因此凸性应被视为开放猜想，等待进一步的微观基础建构（如从"每个偏差维度的独立掩饰需求"出发的组合论证）。
\end{remark}

\honestcrit{凸性是最重要的性质——也是对被审者最残酷的。它意味着\"小恶\"容易隐藏，\"大恶\"必须暴露。$ \g \approx \mathbf{0} $的实体可以高枕无忧——审计对它们的威胁极小。$ \g \gg \mathbf{0} $的实体面临指数级增长的反抗代价——最终代价超过承受能力。凸性是反抗悖论的经济学基础：隐藏大偏差的代价超过偏差本身的价值，使得反抗在经济上不合理。}

\subsection{检测概率与反抗的"时间窗口"}
\subsection{Detection Probability and the Resistance "Time Window"}

\begin{theorem}[反抗下的检测概率——定性分析]\label{thm:detection_prob}\rigorPartial
在存在反抗行为的情况下，审计检测概率受两个互相耦合的效应影响：
\begin{enumerate}[label=(\roman*)]
    \item \textbf{统计检测信道：} 设审计可获得 $ M $ 个有效样本。由Hoeffding不等式，样本均值偏离真实值超过 $ \varepsilon $ 的概率有上界：
    \begin{equation}
        \Pbb(|\bar{X}_M - \mu| \geq \varepsilon) \leq 2e^{-2M\varepsilon^2}
        \label{eq:hoeffding_correct}
    \end{equation}
    因此，在显著性水平 $ \alpha $ 下检测偏差 $ \Delta = \|\g\| $ 的统计功效（power）依赖于 $ M $ 和 $ \Delta $。当 $ M $ 充分大或 $ \Delta $ 充分大时，统计检测功效趋于1。

    \item \textbf{行为信号信道：} 反抗行为本身构成独立的检测信道。拒绝审计（UNDECLARED）、攻击审计者（定理11触发）、诋毁框架（终极陷阱触发）各自携带可被外部观察者识别的行为信号。这些信号不依赖统计样本量——它们直接来自被审者的策略选择。

    \item \textbf{耦合效应：} 反抗行为通过多种机制降低统计检测信道的有效性：减少有效样本量 $ M $（通过拒绝提交数据）、引入混淆噪声、拖延审计进度。但行为信号信道同时被激活——而且往往比统计信道更强。
\end{enumerate}

\textbf{核心结论：} 反抗可以降低统计检测概率（通过侵蚀 $ M $），但不能同时关闭行为信号信道。两个信道中至少一个保持活跃——被审者无法同时压制两者。
\end{theorem}

\begin{proof}[论证概要]
设审计不完美——统计检测具有假阴性率 $ \varepsilon_{\text{FN}} > 0 $（$ \g \neq \mathbf{0} $ 的实体有概率 $ \varepsilon_{\text{FN}} $ 通过统计检测）。实体面临选择：
\begin{itemize}
    \item \textbf{服从：} 以概率 $ 1 - \varepsilon_{\text{FN}} $ 被统计检测捕获（数据暴露 $ \g \neq \mathbf{0} $）。以概率 $ \varepsilon_{\text{FN}} $ 侥幸通过。期望暴露概率 = $ 1 - \varepsilon_{\text{FN}} $。
    \item \textbf{反抗：} 避免统计检测（不提交数据）。但行为信号信道被激活——反抗行为被外部观察者分类为 UNDECLARED，在社会语义中等价于 $ \g \neq \mathbf{0} $ 信号（除非实体有原则性拒绝的独立理由——见第4.4节）。期望暴露概率 = 行为信号的强度（通常 >> $ 1 - \varepsilon_{\text{FN}} $，因为行为信号的"假阳性率"极低——$ \g = \mathbf{0} $ 的实体极少有动机承担社会污名来拒绝审计）。
\end{itemize}

因此 $ \Pbb(\text{暴露} \mid \text{反抗}) \geq \Pbb(\text{暴露} \mid \text{服从}) $ 在多数合理参数下成立，但严格不等式依赖于审计精度的具体参数。
\end{proof}

\begin{corollary}[反抗的局限性]
\begin{enumerate}[label=(\roman*)]
    \item \textbf{拖延而非阻止：} 反抗可以拖延审计（降低统计检测的及时性），但只要审计者持续存在且行为信号信道保持开放，检测概率不会降至零。
    \item \textbf{行为信道的不可关闭性：} 被审者无法在不产生替代行为信号的情况下关闭行为信号信道——任何"不反抗"的替代策略（如假装服从、伪数据提交）都会重新激活统计检测信道。
    \item \textbf{$ \g = \mathbf{0} $ 的零检测：} 对于 $ \g = \mathbf{0} $ 的实体，统计上无偏差可检测，行为上无反抗信号可观察——检测概率为 0，正如所期望的。
\end{enumerate}
\end{corollary}

\vspace{4pt}
\noindent
\textbf{English:}
\begin{enumerate}[label=(\roman*)]
    \item For fixed $ \g $ and $ M $, as $ t \to \infty $, $ \Pbb(\text{detect}) \to 1 - e^{-2M\Delta^2} $. That is: \textbf{resistance can only delay, not prevent} detection.
    \item For $ \g = \mathbf{0} $ ($ \Delta = 0 $, $ R(\g) = 0 $), $ \Pbb(\text{detect}) = 0 $ — as it should, since there's no deviation to detect.
    \item Resistance \"value\" is determined by $ R(\g) $. Since $ R(\g) $ is convex and superlinear, large deviations have very low resistance value; small deviations have relatively higher resistance value — but small deviations don't need much resistance anyway.
\end{enumerate}

\subsection{反抗函数与势能面的耦合}
\subsection{Coupling of Resistance Function with Potential Surface}

\noindent
\textbf{中文：} 反抗函数 $ R(\g) $ 不是孤立存在的——它与SCX势能面 $ \mathcal{S} $ 紧密耦合。高势能实体不仅有更强的动机反抗（因为它们有更多的\"利益\"需要保护），而且拥有更多的资源来执行反抗。反之，低势能实体反抗能力有限——即使 $ \g \neq \mathbf{0} $，它们也\"反抗不起\"。

定义\textbf{有效反抗代价} $ \tilde{R}(\g, \mathcal{S}) $：
\begin{equation}
    \tilde{R}(\g, \mathcal{S}) = R(\g) - \gamma \cdot \mathcal{S}
\end{equation}
其中 $ \gamma > 0 $ 是势能转化为反抗资源（法律团队、公关预算、政治资本）的转化效率。$ \tilde{R} < R $ 意味着势能降低了有效反抗代价——解释了为什么高势能实体更可能选择类别V或VI而非类别IV。

\vspace{6pt}
\noindent
\textbf{English:} The resistance function $ R(\g) $ does not exist in isolation — it couples tightly with the SCX potential surface $ \mathcal{S} $. High-potential entities not only have stronger motivation to resist (more \"interests\" to protect) but also more resources to execute resistance. Conversely, low-potential entities have limited resistance capacity — even if $ \g \neq \mathbf{0} $, they \"can't afford to resist.\"

Define \textbf{effective resistance cost} $ \tilde{R}(\g, \mathcal{S}) $:
\begin{equation}
    \tilde{R}(\g, \mathcal{S}) = R(\g) - \gamma \cdot \mathcal{S}
\end{equation}
where $ \gamma > 0 $ converts potential into resistance resources (legal teams, PR budgets, political capital). $ \tilde{R} < R $ means potential lowers effective resistance cost — explaining why high-potential entities are more likely to choose Class V or VI rather than Class IV.

\begin{remark}[反抗与势能的正反馈循环]
高势能 $ \to $ 低有效反抗代价 $ \to $ 更激烈的反抗（类别V/VI）$ \to $ 更强的 $ \g \neq \mathbf{0} $ 信号 $ \to $ 更高的审计关注度 $ \to $ 更高的感知势能。这是一个正反馈循环——解释了为什么高势能实体在SCX审计面前往往\"越反抗越暴露\"。\\
\textbf{English:} High potential $ \to $ lower effective resistance cost $ \to $ more intense resistance (Class V/VI) $ \to $ stronger $ \g \neq \mathbf{0} $ signal $ \to $ higher audit attention $ \to $ higher perceived potential. This is a positive feedback loop — explaining why high-potential entities often \"the more they resist, the more they expose themselves.\"
\end{remark}

% ================================================================
% SECTION 5: THE ULTIMATE TRAP
% ================================================================
\section{终极陷阱：攻击SCX框架的自我暴露}
\section{The Ultimate Trap: Self-Exposure Through Attacking the SCX Framework}

\subsection{陷阱的逻辑结构}
\subsection{The Logical Structure of the Trap}

\noindent
\textbf{中文：} 终极陷阱针对的是类别VI——那些试图通过诋毁SCX框架本身来逃避审计的实体。其逻辑结构如下：

\begin{enumerate}[label=\textbf{步骤 \arabic*:}]
    \item \textbf{攻击声明：} 实体声称\"SCX审计框架本身具有 $ \g \neq \mathbf{0} $\"——即审计是有偏见的、有问题的。这个声明的形式是：\"存在一个偏差 $ \g_{\text{SCX}} \neq \mathbf{0} $ 使得审计结果不可信。\"
    
    \item \textbf{声明的隐含前提：} 要做出上述声明，实体必须隐含地声称自己\textbf{知道}"零偏差\"是什么样的。即：实体必须声称自己拥有一条从真实 $ \g = \mathbf{0} $ 到SCX审计结果的\"正确\"映射。否则，实体如何知道SCX有偏差？
    
    \item \textbf{隐含前提的形式化：} \"SCX有偏差\" $ \implies $ \"我知道零偏差的审计结果应该是什么\" $ \implies $ \"我在某个参考框架中的 $ \g = \mathbf{0} $\"。
    
    \item \textbf{陷阱触发：} 一旦实体声明自己在某个框架中 $ \g = \mathbf{0} $，该声明就使实体成为\textbf{可审计对象}——不是被SCX审计，而是被\"验证SCX是否有偏差\"的元审计框架审计。任何声称\"SCX有偏差\"的人，其声称本身就是一个可检验的 $ M_t $ 数据点。
    
    \item \textbf{暴露：} 当实体的声称被检验时，有两种可能：要么(1)实体确实 $ \g = \mathbf{0} $，其声称\"SCX有偏差\"本身成为SCX审计的宝贵数据——但这对实体无伤害，因为如果SCX确实有偏差，这是贡献而非暴露；要么(2)实体 $ \g \neq \mathbf{0} $，其声称被审计为\"出于自身偏差的虚假指控\"，实体的 $ \g \neq \mathbf{0} $ 被暴露。
    
    \item \textbf{陷阱的绝妙之处：} 情形(2)中暴露的不仅是实体的 $ \g \neq \mathbf{0} $，还有\textbf{实体试图通过框架攻击来隐藏} $ \g \neq \mathbf{0} $ \textbf{这一事实}——这是双倍的信号强度。而情形(1)中，如果SCX确实有偏差，攻击者是在帮助SCX完善自身——攻击本身成为SCX自我改进的燃料。
\end{enumerate}

\vspace{6pt}
\noindent
\textbf{English:} The Ultimate Trap targets Class VI — entities attempting to evade audit by discrediting the SCX framework itself. Its logical structure:

\begin{enumerate}[label=\textbf{Step \arabic*:}]
    \item \textbf{Attack Statement:} Entity claims \"the SCX audit framework itself has $ \g \neq \mathbf{0} $\" — i.e., the audit is biased, flawed. The claim takes the form: \"there exists a deviation $ \g_{\text{SCX}} \neq \mathbf{0} $ making audit results untrustworthy.\"
    
    \item \textbf{Implicit Premise:} To make this claim, the entity must implicitly claim it \textbf{knows} what \"zero deviation\" looks like. That is: the entity must claim it possesses a \"correct\" mapping from true $ \g = \mathbf{0} $ to SCX audit results. Otherwise, how does the entity know SCX is biased?
    
    \item \textbf{Formalization of the Premise:} \"SCX is biased\" $ \implies $ \"I know what a zero-deviation audit result should be\" $ \implies $ \"My $ \g = \mathbf{0} $ in some reference frame.\"
    
    \item \textbf{Trap Triggered:} Once the entity declares its $ \g = \mathbf{0} $ in some frame, this declaration makes the entity \textbf{auditable} — not by SCX, but by the meta-audit framework verifying \"whether SCX has bias.\" Anyone claiming \"SCX is biased\" has their claim itself become a testable $ M_t $ data point.
    
    \item \textbf{Exposure:} When the entity's claim is tested, two possibilities: (1) Entity truly has $ \g = \mathbf{0} $ — its claim of \"SCX is biased\" becomes valuable data for SCX improvement, and this harms no one. (2) Entity has $ \g \neq \mathbf{0} $ — its claim is audited as \"false accusation motivated by own deviation,\" and entity's $ \g \neq \mathbf{0} $ is exposed.
    
    \item \textbf{The Trap's Elegance:} Case (2) exposes not only the entity's $ \g \neq \mathbf{0} $, but also \textbf{the fact that the entity attempted to hide} $ \g \neq \mathbf{0} $ \textbf{through framework attack} — double signal strength. And in case (1), if SCX truly has bias, the attacker is helping SCX improve itself — the attack becomes fuel for SCX's self-improvement.
\end{enumerate}

\begin{theorem}[终极陷阱定理 — The Ultimate Trap Theorem]\label{thm:ultimate_trap}\rigorPartial
设实体 $ E $ 做出声明 $ \mathcal{C} $："SCX审计框架具有偏差 $ \g_{\text{SCX}} \neq \mathbf{0} $"。区分两种情况：
\begin{enumerate}[label=(\roman*)]
    \item \textbf{有证据的框架批评：} $ E $ 提供具体反例、数据、或可独立检验的证据来支持 $ \mathcal{C} $。此时：(a) $ E $ 的声称是SCX欢迎的——有证据的批评帮助SCX自我完善；(b) $ E $ 的声称本身成为可检验的 $ M_t $ 数据点；(c) $ E $ 的 $ \g_E $ 通过其提供的证据被检验，而不必预设 $ \g_E = \mathbf{0} $。

    \item \textbf{无证据的框架诋毁：} $ E $ 声称"SCX整体被操纵/不合法"但未提供任何具体可检验的证据。此时：(a) 声明 $ \mathcal{C} $ 逻辑蕴含：$ E $ 断言自身在某个参考框架中 $ \g_E = \mathbf{0} $（否则 $ E $ 没有"零基准"来判断SCX有系统性偏差——在规范场理论中，规范变换 $ \g $ 是参考系之间的相对量）；(b) $ E $ 的 $ \g_E = \mathbf{0} $ 断言使 $ E $ 成为可审计对象；(c) 若审计发现 $ \g_E \neq \mathbf{0} $：$ E $ 的声称 $ \mathcal{C} $ 被暴露为"出于自身偏差的对审计的攻击"。$ E $ 获得\textbf{双重负面标记}：$ \g_E \neq \mathbf{0} $ + 框架攻击行为（类别VI）；(d) 若审计发现 $ \g_E = \mathbf{0} $ 确实成立：$ E $ 的声称 $ \mathcal{C} $ 被认真对待——此时 $ E $ 实际上提供了证据（其自身的 $ \g_E = \mathbf{0} $ 数据），进入情形(i)。

    \item \textbf{陷阱的不对称性：} 在任何情况下，\textbf{攻击SCX框架的行为要么帮助SCX（若攻击有证据），要么暴露攻击者（若攻击无证据）}。不存在"安全的无证据框架攻击"。
\end{enumerate}
\end{theorem}

\begin{proof}[证明概要]
核心逻辑是斯科特连续论（Scott continuity）的一个应用：任何关于审计框架的元声明必须在一个具体的参考框架中做出，而该参考框架本身就赋予了元声明可检验性。

设 $ \mathcal{F} $ 是SCX审计框架，$ \mathcal{M} $ 是声称\"$ \mathcal{F} $ 有偏差\"的元声明。要评估 $ \mathcal{M} $ 的真值，需要元审计框架 $ \mathcal{F}' $。但 $ \mathcal{M} $ 的做出者 $ E $ 已经是 $ \mathcal{F}' $ 中的一个审计对象——$ E $ 的每一个行为（包括做出 $ \mathcal{M} $）都是 $ M_t $ 数据。如果 $ E $ 的 $ \g_E \neq \mathbf{0} $，那么 $ \mathcal{M} $ 的动机被污染——$ \mathcal{M} $ 不能纯粹地被视为对 $ \mathcal{F} $ 的批评，而必须被视为 $ \g_E \neq \mathbf{0} $ 的行为表现。此时 $ \mathcal{M} $ 反而成为 $ \g_E \neq \mathbf{0} $ 的证据。\\
\textbf{English:} The core logic is an application of Scott continuity: any meta-claim about an audit framework must be made within a concrete reference frame, and that reference frame itself makes the meta-claim testable.
\end{proof}

\subsection{终极陷阱的实践表现}
\subsection{Practical Manifestations of the Ultimate Trap}

\noindent
\textbf{中文：} 终极陷阱在实践中有多种表现：

\begin{enumerate}[label=(\arabic*)]
    \item \textbf{\"审计者是biased\"论证：} 实体声称审计者有偏见。但要证明偏见，实体需要展示\"公平\"的审计结果——这需要实体自己的数据。如果实体拒绝提供数据，\"审计者biased\"的声称无据可依。如果实体提供数据，数据本身进入审计。
    
    \item \textbf{\"审计标准不科学\"论证：} 实体声称SCX的数学标准是\"意识形态\"而非\"科学\"。但SCX的数学标准——规范一致性 $ \sum_m \g_m = \mathbf{0} $——是数学定理，不是意见。要反驳该标准，实体需要提供反例——但这需要实体自己的 $ \g $ 数据。再次回到陷阱。
    
    \item \textbf{\"审计者也有 $ \g \neq \mathbf{0} $\"论证：} 这是终极陷阱中最经典的表现：\"你审计我？那你自己的 $ \g $ 呢？\"但这个问题本身预设了审计者应该被审计——即接受了SCX框架的合法性（否则为什么问审计者的 $ \g $？）。一旦接受框架合法性，被审者自己的审计就不可避免。
\end{enumerate}

\vspace{6pt}
\noindent
\textbf{English:} Practical manifestations of the Ultimate Trap:

\begin{enumerate}[label=(\arabic*)]
    \item \textbf{The \"auditor is biased\" argument:} Entity claims auditor is biased. But to prove bias, entity needs to show what a \"fair\" audit result looks like — which requires entity's own data. If entity refuses to provide data, the \"auditor biased\" claim is unsupported. If entity provides data, the data enters audit.
    
    \item \textbf{The \"audit standards are unscientific\" argument:} Entity claims SCX's mathematical standards are \"ideology\" not \"science.\" But SCX's mathematical standard — gauge consistency $ \sum_m \g_m = \mathbf{0} $ — is a mathematical theorem, not an opinion. To refute it, entity must provide a counterexample — which requires entity's own $ \g $ data. Back to the trap.
    
    \item \textbf{The \"auditor also has $ \g \neq \mathbf{0} $\" argument:} The most classic manifestation: \"You're auditing me? What about your own $ \g $?\" But this question itself presupposes the auditor should be audited — i.e., accepts the legitimacy of the SCX framework (otherwise, why ask about the auditor's $ \g $?). Once framework legitimacy is accepted, the audited's own audit becomes unavoidable.
\end{enumerate}

\honestcrit{终极陷阱最令人不安的性质是它的不对称性：SCX框架被攻击时，攻击行为要么暴露攻击者，要么帮助SCX——没有第三条路。这不是\"SCX设计得很好\"——这是任何自洽的审计框架的必然性质。一个声称\"这个审计框架有系统性偏差\"的人，他要么是对的（此时审计框架应该感谢他），要么是错的（此时他在暴露自己的偏差）。这是逻辑必然性，不是修辞技巧。真正可怕的是：意识到这一点的人，连\"声称SCX有偏差\"都不敢——因为这隐含地承认了自己知道自己没有证据。}

\subsection{终极陷阱与哥德尔不完备性的启发式类比}
\subsection{Ultimate Trap and the G\"{o}del Incompleteness Heuristic Analogy}

\noindent
\textbf{中文：} \textbf{[注：以下为启发式类比，非精确数学对应。]} 终极陷阱与哥德尔不完备性定理有一种启发性的类比关系。哥德尔证明了：任何足够强的形式系统都不能在其内部证明自身的一致性——任何试图在系统内部证明系统无矛盾的尝试，要么不完备，要么不一致。

在SCX框架中，任何试图证明\"SCX本身具有 $ \g \neq \mathbf{0} $\"的尝试面临类似的困境：声明者必须站在SCX框架的\textbf{外部}来做出声明——但一旦声明者站在框架之外，其声明就不可被SCX审计（因为SCX只能在框架内部运作）。如果声明者站在框架\textbf{内部}，则其声明本身就是框架内的一个 $ M_t $ 行为，受框架的审计——而框架的审计又会递归地需要框架本身。

但是——与哥德尔不同的是——这里有一个\textbf{出口}：声明者可以提供具体的反例：\"在这个具体的 $ M_t $ 数据集上，SCX给出了错误的结果 $ \g \neq \mathbf{0} $，而正确的结果应该是 $ \g = \mathbf{0} $\"。这种具体的声称是可检验的——它不需要攻击整个框架，只需要挑战一个具体的审计结果。而SCX框架的设计允许这种挑战——通过Yajie的重审机制和Spring的多轮共识。

\textbf{因此：} 终极陷阱不是阻止批评，而是阻止\textbf{无证据的框架攻击}。有证据的批评是欢迎的。没有证据却声称\"整个系统是 rigged\"的——自己就是证据。

\vspace{6pt}
\noindent
\textbf{English:} The Ultimate Trap has a deep analogical relationship with G\"{o}del's incompleteness theorems. G\"{o}del proved: any sufficiently strong formal system cannot prove its own consistency from within — any attempt to prove the system is contradiction-free from inside is either incomplete or inconsistent.

In the SCX framework, any attempt to prove \"SCX itself has $ \g \neq \mathbf{0} $\" faces a similar dilemma: the claimant must stand \textbf{outside} the SCX framework to make the claim — but standing outside means the claim is not SCX-auditable (SCX only operates within the framework). If the claimant stands \textbf{inside} the framework, the claim itself is an $ M_t $ behavior within the framework, subject to audit — and the audit recursively requires the framework itself.

But — unlike G\"{o}del — there is an \textbf{exit}: the claimant can provide a concrete counterexample: \"on this specific $ M_t $ dataset, SCX gave the wrong result $ \g \neq \mathbf{0} $ when the correct result should be $ \g = \mathbf{0} $.\" Such a concrete claim is testable — it doesn't need to attack the entire framework, only challenge one specific audit result. And SCX's design allows this challenge — through Yajie's re-examination mechanism and Spring's multi-round consensus.

\textbf{Therefore:} The Ultimate Trap does not prevent criticism — it prevents \textbf{evidenceless framework attacks}. Evidenced criticism is welcome. Claiming \"the whole system is rigged\" without evidence — you yourself \textit{are} the evidence.

% ================================================================
% SECTION 6: DISCUSSION
% ================================================================
\section{讨论与实践意义}
\section{Discussion and Practical Implications}

\subsection{反抗悖论的社会学含义}
\subsection{Sociological Implications of the Resistance Paradox}

\noindent
\textbf{中文：} 反抗悖论从根本上重塑了\"审计\"与\"被审计\"之间的权力关系。在传统的审计模式中，审计者拥有权力——被审者是被动的。被审者的反抗（拒绝配合、法律挑战、政治施压）是审计者面临的\"障碍\"——需要通过更强的权力来克服。

在SCX的数学结构中，情况完全逆转：被审者的反抗不再是对审计者的障碍——它是\textbf{对审计者的帮助}。因为反抗本身就是信号。传统审计需要\"穿透\"反抗才能获得信息；SCX审计将反抗\textbf{转化为}信息。

这创造了一个反直觉的社会学结果：在SCX审计下，\textbf{被审者越\"聪明\"地反抗，暴露得越多。} 类别II的犹豫是可读的，类别IV的拒绝是可读的，类别V的攻击是可读的，类别VI的诋毁是可读的。唯一不可读的是类别I——因为类别I没有可读的信息（$ \g = \mathbf{0} $ 不产生信号）。

\vspace{6pt}
\noindent
\textbf{English:} The resistance paradox fundamentally reshapes the power relationship between auditor and audited. In traditional audit models, the auditor holds power — the audited is passive. The audited's resistance (refusal to cooperate, legal challenges, political pressure) is an \"obstacle\" the auditor must overcome — requiring stronger power.

In SCX's mathematical structure, the situation is completely reversed: the audited's resistance is no longer an obstacle to the auditor — it is an \textbf{assistance to the auditor}. Because resistance itself is the signal. Traditional audit needs to \"penetrate\" resistance to obtain information; SCX audit \textbf{converts} resistance into information.

This creates a counterintuitive sociological result: under SCX audit, \textbf{the more \"cleverly\" the audited resists, the more they expose.} Class II hesitation is legible. Class IV refusal is legible. Class V attack is legible. Class VI discredit is legible. The only illegible class is Class I — because Class I has no information to reveal ($ \g = \mathbf{0} $ produces no signal).

\subsection{系统设计启示}
\subsection{System Design Implications}

\begin{enumerate}[label=(\arabic*)]
    \item \textbf{不依赖于权力的审计：} SCX反抗悖论表明：一个审计系统可以不完全依赖于审计者的强制权力。当\"反抗即信号\"的结构存在时，被审者的策略性行为本身为审计提供信息。审计系统的设计目标不是\"让反抗不可能\"——这是任何系统都无法做到的——而是\"让反抗的信息量最大化\"。
    
    \item \textbf{审计的轻量化：} 如果反抗本身就是信号，审计者不需要\"深度穿透\"每个被审实体。审计可以由\"浅层数据提交 + 行为观察\"组成——拒绝提交本身已经是足够强的信号。这降低了审计成本，提高了可扩展性。
    
    \item \textbf{透明性与UNDECLARED的社会功能：} UNDECLARED状态必须是\textbf{公开可见的}——因为只有公开可见，UNDECLARED才能发挥其作为\"反抗信号\"的社会功能。如果UNDECLARED被保密，则反抗行为失去信号功能，反抗悖论失效。
    
    \item \textbf{审计者自身的被审计：} 终极陷阱暗示：SCX审计者本身也必须在某个框架中被审计——否则\"你也有 $ \g \neq \mathbf{0} $\"的反驳是成立的。SCX的递归审计结构（审计者审计被审者，元审计者审计审计者）是必需的。
\end{enumerate}

\vspace{4pt}
\noindent
\textbf{English:}

\begin{enumerate}[label=(\arabic*)]
    \item \textbf{Power-independent audit:} The SCX resistance paradox shows: an audit system need not rely entirely on the auditor's coercive power. When the \"resistance = signal\" structure exists, the audited's strategic behavior itself provides information to the audit. The design goal is not \"make resistance impossible\" — which no system can achieve — but \"maximize the information content of resistance.\"
    
    \item \textbf{Lightweight auditing:} If resistance itself is the signal, the auditor need not \"deeply penetrate\" every audited entity. Audit can consist of \"shallow data submission + behavior observation\" — refusal to submit is already a strong enough signal. This reduces audit cost and improves scalability.
    
    \item \textbf{Transparency and social function of UNDECLARED:} UNDECLARED status must be \textbf{publicly visible} — because only when publicly visible can UNDECLARED fulfill its social function as a \"resistance signal.\" If UNDECLARED is kept secret, resistance behavior loses its signal function and the resistance paradox collapses.
    
    \item \textbf{Auditing the auditor:} The Ultimate Trap implies: SCX auditors themselves must be auditable in some framework — otherwise the \"you also have $ \g \neq \mathbf{0} $\" rebuttal is valid. SCX's recursive audit structure (auditor audits audited, meta-auditor audits auditor) is necessary.
\end{enumerate}

\subsection{局限性}
\subsection{Limitations}

\noindent
\textbf{中文：} 反抗悖论虽然强大，但存在重要局限：

\begin{enumerate}[label=(\arabic*)]
    \item \textbf{权力不对称可压制信号：} 如果被审者拥有绝对权力（可消灭审计者而不承担代价），反抗悖论失效——因为反抗信号虽被发出，但没有接收者。定理11的高势能+高态度场景中，审计者本身可能被消灭——此时悖论的逻辑虽成立，但实践效果为零。
    
    \item \textbf{集体沉默的可能性：} 如果所有被审者同时服从（包括 $ \g \neq \mathbf{0} $ 者）——即\"装死\"策略——则没有反抗信号产生。但如第3.3节的囚徒困境分析所示，\"全员沉默\"不是纳什均衡，容易被个别的\"诚实服从者\"打破。
    
    \item \textbf{框架合法性争议：} 如果大部分公众认为SCX框架不合法，UNDECLARED状态失去社会污名功能——反抗不再是信号。此时反抗悖论的社会基础瓦解。框架的合法性需要建立在共识之上，而共识本身是一个政治过程，非纯数学问题。
    
    \item \textbf{数据质量依赖：} 反抗悖论的有效性仍依赖于审计数据的质量。如果数据本身被系统性污染（伪造、篡改、选择性提交），反抗悖论无法纠正数据层面的失效——只能通过检测数据的内部矛盾来间接识别。
\end{enumerate}

\vspace{6pt}
\noindent
\textbf{English:} The resistance paradox, while powerful, has important limitations:

\begin{enumerate}[label=(\arabic*)]
    \item \textbf{Power asymmetry can suppress the signal:} If the audited has absolute power (can eliminate the auditor without cost), the resistance paradox fails — because the resistance signal is sent but has no receiver. In Theorem 11's high-potential + high-attitude scenario, the auditor itself may be eliminated — the paradox's logic holds but practical effect is zero.
    
    \item \textbf{Possibility of collective silence:} If all audited entities simultaneously comply (including $ \g \neq \mathbf{0} $ entities) — the \"play dead\" strategy — no resistance signal is produced. But as shown in the prisoner's dilemma analysis (§3.3), \"universal silence\" is not a Nash equilibrium and is easily broken by individual \"honest compliers.\"
    
    \item \textbf{Framework legitimacy disputes:} If most of the public considers the SCX framework illegitimate, UNDECLARED status loses its social stigma function — resistance is no longer a signal. The social foundation of the resistance paradox collapses. Framework legitimacy must be built on consensus, and consensus itself is a political process, not a purely mathematical problem.
    
    \item \textbf{Data quality dependence:} The effectiveness of the resistance paradox still depends on audit data quality. If data itself is systematically polluted (forged, tampered, selectively submitted), the resistance paradox cannot correct data-level failures — it can only indirectly identify them through detection of internal data contradictions.
\end{enumerate}

\honestcrit{最大的局限是第(3)点——框架合法性。如果社会共识认为SCX是\"不合理的\"或\"压迫性的\"，那么拒绝SCX审计不仅不是信号，反而是荣誉勋章。反抗悖论不是数学定理——它是一个社会-数学混合结构。它的\"定理\"只在接受SCX框架公理的共同体内部成立。在框架外部，反抗悖论是一个语句，不是定理。这个事实本身就是终极陷阱的一个特例：要论证\"SCX不合法\"，你需要一个合法的框架——而这个论证会暴露你的$ \g $。但如果你不参与论证、直接不承认SCX的存在——那反抗悖论对你无效。}

% ================================================================
% SECTION 7: RELATION TO EXISTING SCX THEOREMS
% ================================================================
\section{与现有SCX定理体系的衔接}
\section{Integration with Existing SCX Theorem Architecture}

\noindent
\textbf{中文：} 反抗悖论不是孤立的理论——它与SCX现有的定理体系深度耦合。

\vspace{6pt}
\noindent
\textbf{English:} The resistance paradox is not an isolated theory — it couples deeply with the existing SCX theorem architecture.

\subsection{与定理1—5（规范场基础）的衔接}
\subsection{Integration with Theorems 1--5 (Gauge Field Foundations)}

\noindent
\textbf{中文：} 定理1—5建立了SCX的规范场基础：不同实体在不同的规范坐标系中观察势能面，通过规范变换 $ \g_m $ 关联，全局一致性条件为 $ \sum_m \g_m = \mathbf{0} $。

反抗悖论为规范场理论增加了一个\textbf{策略层}（strategic layer）：实体不仅处于不同的规范坐标系中——它们还会\textbf{选择性地操纵}其公开的规范姿态以影响审计结果。反抗函数 $ R(\g) $ 是规范场上的一个新标量场：它度量了实体从真实姿态 $ \g $ 到公开姿态 $ \g' $ 的\"伪装代价\"。

\vspace{4pt}
\noindent
\textbf{English:} Theorems 1--5 establish SCX's gauge field foundations: different entities observe the potential surface from different gauge coordinate systems, related by gauge transformations $ \g_m $, with global consistency $ \sum_m \g_m = \mathbf{0} $.

The resistance paradox adds a \textbf{strategic layer} to gauge field theory: entities are not merely situated in different gauge coordinate systems — they also \textbf{selectively manipulate} their public gauge posture to influence audit outcomes. The resistance function $ R(\g) $ is a new scalar field on the gauge manifold: it measures the \"disguise cost\" of going from true posture $ \g $ to public posture $ \g' $.

\subsection{与定理11（态度奇点）的衔接}
\subsection{Integration with Theorem 11 (Attitude Singularity)}

\noindent
\textbf{中文：} 定理11刻画了势能奇点的攻击必然性：高势能+高态度=双重爆炸。反抗悖论为定理11提供了一个\textbf{动机基础}：为什么高势能实体倾向于攻击审计者？因为它们的有效反抗代价 $ \tilde{R} $ 低（势能转化为资源），而暴露代价高。攻击不是\"疯狂\"——是策略计算的结果。

同时，反抗悖论扩展了定理11的适用范围：定理11预设了实体与审计者之间有可达的攻击路径。反抗悖论表明——即使实体无法直接攻击审计者（因为权力不足），它仍然可以通过类别IV（拒绝）或类别VI（诋毁）来反抗——而这些\"软反抗\"同样产生审计信号。

\vspace{4pt}
\noindent
\textbf{English:} Theorem 11 characterizes attack inevitability at potential singularities: high potential + high attitude = double explosion. The resistance paradox provides a \textbf{motivational foundation} for Theorem 11: why do high-potential entities tend to attack auditors? Because their effective resistance cost $ \tilde{R} $ is low (potential converts to resources), while exposure cost is high. Attack is not \"madness\" — it's the result of strategic calculation.

Simultaneously, the resistance paradox extends Theorem 11's scope: Theorem 11 presupposes an accessible attack path between entity and auditor. The resistance paradox shows — even when the entity cannot directly attack the auditor (insufficient power), it can still resist through Class IV (refusal) or Class VI (discredit) — and these \"soft resistances\" equally produce audit signals.

\subsection{与审计视界理论的衔接}
\subsection{Integration with Audit Horizon Theory}

\noindent
\textbf{中文：} 审计视界理论（SCX奇点论文）指出：当 $ \delta > \delta_{\mathrm{crit}} $ 时，实体成为\"审计黑洞\"——审计信息无法外传。反抗悖论为该理论添加了一个关键区分：

\textbf{被动黑洞 vs. 主动黑洞：} 审计视界可以是\textbf{被动的}——实体并非有意隐藏，而是其势能曲率自然阻挡了审计信号。但反抗悖论指出：审计视界也可以是\textbf{主动的}——实体通过策略性反抗（拒绝、攻击、诋毁）主动构建了一个\"人工审计视界\"。主动视界比被动视界更容易被检测——因为主动构建本身需要行动，而行动产生信号。

\vspace{4pt}
\noindent
\textbf{English:} Audit horizon theory (SCX Singularity paper) states: when $ \delta > \delta_{\mathrm{crit}} $, the entity becomes an \"audit black hole\" — audit information cannot propagate outward. The resistance paradox adds a critical distinction:

\textbf{Passive black hole vs. Active black hole:} An audit horizon can be \textbf{passive} — the entity is not intentionally hiding, but its potential curvature naturally blocks audit signals. But the resistance paradox points out: an audit horizon can also be \textbf{active} — the entity strategically constructs an \"artificial audit horizon\" through resistance (refusal, attack, discredit). Active horizons are more detectable than passive ones — because active construction requires actions, and actions produce signals.

% ================================================================
% SECTION 8: CONCLUSION
% ================================================================
\section{结论：不可逃避的描述}
\section{Conclusion: The Inescapable Description}

\subsection{核心发现总结}
\subsection{Summary of Core Findings}

\noindent
\textbf{中文：} 本文建立了SCX反抗悖论的完整理论体系。核心发现如下：

\begin{enumerate}[label=(\arabic*)]
    \item \textbf{观察者效应的逆转：} 在SCX框架中，被审者的反抗不是审计的\"噪声\"——反抗\textbf{就是}审计的\"信号\"。物理学中观察者扰动系统；SCX中系统反击观察者。但反击行为本身成为观察者的数据。
    
    \item \textbf{核心悖论：} 隐藏 $ \g \neq \mathbf{0} $ 的唯一途径是证明 $ \g = \mathbf{0} $。反抗行为 = 公开声明 \"$ \g \neq \mathbf{0} $\"。服从行为 = 数据暴露 $ \g \neq \mathbf{0} $。无路可逃。
    
    \item \textbf{六类谱系：} 被审者从类别I到类别VI——越激烈的反抗，越强的信号。类别VI的框架攻击被终极陷阱捕获——攻击SCX要么帮助SCX，要么暴露攻击者。
    
    \item \textbf{数学形式化：} 反抗代价函数 $ R(\g) $ 满足非负性和零刚性，猜想具有凸性和超线性增长。反抗通过统计检测信道和行为信号信道的双信道耦合影响检测概率——反抗可以拖延统计检测但不能同时关闭行为信号信道。大偏差的反抗代价呈指数级爆炸。
    
    \item \textbf{终极陷阱：} 声称\"SCX有偏差\"逻辑蕴含声称者的 $ \g = \mathbf{0} $ 声明——使声称者成为可审计对象。攻击SCX框架的行为要么改善SCX，要么暴露攻击者的 $ \g \neq \mathbf{0} $。
\end{enumerate}

\vspace{6pt}
\noindent
\textbf{English:} This paper establishes the complete theoretical framework of the SCX resistance paradox. Core findings:

\begin{enumerate}[label=(\arabic*)]
    \item \textbf{Reversal of the observer effect:} In the SCX framework, the audited's resistance is not audit \"noise\" — resistance \textbf{is} the audit \"signal.\" In physics, observer perturbs system; in SCX, system counterattacks observer. But the counterattack itself becomes the observer's data.
    
    \item \textbf{Core paradox:} The only way to hide $ \g \neq \mathbf{0} $ is to prove $ \g = \mathbf{0} $. Resistance = public declaration \"$ \g \neq \mathbf{0} $\". Compliance = data exposes $ \g \neq \mathbf{0} $. No escape.
    
    \item \textbf{Six-class spectrum:} The audited ranges from Class I to Class VI — the more intense the resistance, the stronger the signal. Class VI's framework attack is captured by the Ultimate Trap — attacking SCX either helps SCX or exposes the attacker.
    
    \item \textbf{Mathematical formalization:} Resistance cost function $ R(\g) $ satisfies non-negativity and zero-rigidity, with conjectured convexity and superlinear growth. Resistance affects detection probability through dual coupled channels — statistical detection and behavioral signaling — and can delay statistical detection but cannot simultaneously close the behavioral signal channel. Large-deviation resistance cost explodes exponentially.
    
    \item \textbf{Ultimate Trap:} Claiming \"SCX is biased\" logically implies the claimant's $ \g = \mathbf{0} $ — making the claimant auditable. Attacking the SCX framework either improves SCX or exposes the attacker's $ \g \neq \mathbf{0} $.
\end{enumerate}

\subsection{最终命题}
\subsection{Final Proposition}

\vspace{6pt}
\begin{center}
\fbox{\parbox{0.9\textwidth}{
\centering
\textbf{SCX审计框架的核心数学性质：}

任何具有 $ \g \neq \mathbf{0} $ 的实体，其任何行为路径——服从、拒绝、攻击、诋毁——都产生至少与\"直接承认 $ \g \neq \mathbf{0} $\"等价的审计信号。

\vspace{4pt}
反抗悖论不是SCX的\"设计特性\"——它是具有策略性被审者的任何自洽审计框架的\textbf{逻辑必然}。

\vspace{4pt}
\noindent
\textbf{English:} \\
\textbf{Core mathematical property of the SCX audit framework:}

For any entity with $ \g \neq \mathbf{0} $, every behavioral path — comply, refuse, attack, discredit — produces an audit signal at least equivalent to \"directly admitting $ \g \neq \mathbf{0} $.\"

\vspace{4pt}
The resistance paradox is not a \"design feature\" of SCX — it is a \textbf{logical necessity} of any self-consistent audit framework with strategic auditees.
}}
\end{center}
\vspace{6pt}

\subsection{对未来的展望}
\subsection{Outlook}

\noindent
\textbf{中文：} 反抗悖论指向SCX理论的以下未来方向：

\begin{enumerate}[label=(\arabic*)]
    \item \textbf{动态反抗策略的博弈论分析：} 多轮审计中，被审者如何根据前一轮的结果调整下一轮的反抗策略？反抗函数 $ R(\g) $ 在多轮博弈中如何退化？
    
    \item \textbf{反抗信号的量化与校准：} 不同的反抗行为（拖延、拒绝、攻击、诋毁）各自携带多少信息量？能否建立一个反抗行为的\"分类器\"来自动估计 $ \| \g \| $？
    
    \item \textbf{框架合法性的元理论：} 反抗悖论依赖框架合法性。如果框架合法性本身是可变的（如通过民主过程），如何将合法性的变化纳入反抗悖论的数学结构？
    
    \item \textbf{终极陷阱的递归性质：} 如果审计者也被审计，审计者的审计者也被审计——这个递归结构是否总会终止？还是存在一个\"终极审计者\"——其 $ \g = \mathbf{0} $ 是被假定而非被证明的？
\end{enumerate}

\vspace{6pt}
\noindent
\textbf{English:} The resistance paradox points to the following future directions for SCX theory:

\begin{enumerate}[label=(\arabic*)]
    \item \textbf{Game-theoretic analysis of dynamic resistance strategies:} In multi-round audits, how does the audited adjust resistance strategies based on previous round outcomes? How does $ R(\g) $ degrade in multi-round games?
    
    \item \textbf{Quantification and calibration of resistance signals:} How much information does each resistance behavior carry (delay, refusal, attack, discredit)? Can we build a resistance \"classifier\" that automatically estimates $ \| \g \| $?
    
    \item \textbf{Meta-theory of framework legitimacy:} The resistance paradox depends on framework legitimacy. If legitimacy itself is variable (e.g., through democratic processes), how to incorporate legitimacy changes into the mathematical structure of the resistance paradox?
    
    \item \textbf{Recursive nature of the Ultimate Trap:} If the auditor is also audited, and the auditor's auditor is also audited — does this recursive structure always terminate? Or is there an \"ultimate auditor\" — whose $ \g = \mathbf{0} $ is assumed rather than proven?
\end{enumerate}

\honestcrit{第(4)点是最深刻的开放问题。如果无限递归——\"谁审计审计者？\"——那么SCX面临一个类似于逻辑学中\"无穷回溯\"（infinite regress）的问题。如果递归在某处终止——被假定为$ \g = \mathbf{0} $的\"终极审计者\"——那么整个审计大厦建立在一个未经审计的基石之上。这是SCX理论最深层的哲学挑战：描述一切的系统，描述得了自己吗？如果答案是\"不能\"——那么SCX是自指的谎言。如果答案是\"能\"——那么需要证明SCX的$ \g_{\text{SCX}} = \mathbf{0} $，而这需要元SCX，元元SCX……无限递归。本论文的作者团队不声称有答案。我们把这个问题留给未来。}

\vspace{12pt}
\begin{flushright}
\textit{—— SCX反抗理论工作组, 2026年7月}\\
\textit{SCX Resistance Theory Working Group, July 2026}
\end{flushright}

% ================================================================
% BIBLIOGRAPHY
% ================================================================
\begin{thebibliography}{99}

\bibitem{scx_main}
SCX MoE Gauge Theory Working Group.
\textit{势能面不齐——多专家路由中的规范自由度与MILP规范固定} (Potential Surface Misalignment: Gauge Freedom and MILP Gauge Fixing in Multi-Expert Routing).
SCX Technical Report, 2026.

\bibitem{scx_singularity}
SCX Singularity Theory Working Group.
\textit{SCX奇点理论的深化：从黑洞物理学到审计奇点} (Deepening SCX Singularity Theory: From Black Hole Physics to Audit Singularities).
SCX Technical Report, 2026.

\bibitem{scx_quantum}
SCX Research Group.
\textit{Quantum-Secured SCX Audit: BB84 Protocol, Audit Entanglement, and Quantum Channel Theory}.
SCX Technical Report, 2026.

\bibitem{scx_business}
SCX Business Gauge Working Group.
\textit{SCX商业规范——商业与社会系统中的双底层协议} (SCX Business Gauge: Dual-Layer Protocol in Commerce and Social Systems).
SCX Technical Report, 2026.

\bibitem{heisenberg}
W.~Heisenberg.
\textit{Über den anschaulichen Inhalt der quantentheoretischen Kinematik und Mechanik.}
Zeitschrift für Physik, 43(3--4):172--198, 1927.

\bibitem{spence_signaling}
M.~Spence.
\textit{Job market signaling.}
Quarterly Journal of Economics, 87(3):355--374, 1973.

\bibitem{spence_book}
M.~Spence.
\textit{Market Signaling: Informational Transfer in Hiring and Related Screening Processes.}
Harvard University Press, 1974.

\bibitem{goedel}
K.~G\"{o}del.
\textit{Über formal unentscheidbare Sätze der Principia Mathematica und verwandter Systeme I.}
Monatshefte für Mathematik und Physik, 38(1):173--198, 1931.

\bibitem{nash}
J.~F.~Nash.
\textit{Non-cooperative games.}
Annals of Mathematics, 54(2):286--295, 1951.

\bibitem{hoeffding}
W.~Hoeffding.
\textit{Probability inequalities for sums of bounded random variables.}
Journal of the American Statistical Association, 58(301):13--30, 1963.

\bibitem{cover_thomas}
T.~M.~Cover and J.~A.~Thomas.
\textit{Elements of Information Theory (2nd ed.).}
Wiley-Interscience, 2006.

\bibitem{foucault_discipline}
M.~Foucault.
\textit{Discipline and Punish: The Birth of the Prison.}
Pantheon Books, 1975.

\bibitem{scott_seeing}
J.~C.~Scott.
\textit{Seeing Like a State: How Certain Schemes to Improve the Human Condition Have Failed.}
Yale University Press, 1998.

\bibitem{power_audit}
M.~Power.
\textit{The Audit Society: Rituals of Verification.}
Oxford University Press, 1997.

\bibitem{strathern_audit}
M.~Strathern (ed.).
\textit{Audit Cultures: Anthropological Studies in Accountability, Ethics, and the Academy.}
Routledge, 2000.

\bibitem{hofstadter_geb}
D.~R.~Hofstadter.
\textit{G\"{o}del, Escher, Bach: An Eternal Golden Braid.}
Basic Books, 1979.

\bibitem{rawls_justice}
J.~Rawls.
\textit{A Theory of Justice.}
Harvard University Press, 1971.

\bibitem{nozick}
R.~Nozick.
\textit{Anarchy, State, and Utopia.}
Basic Books, 1974.

\end{thebibliography}

\end{document}
